\section{Lockean threshold phase transition under probability integrity}
\label{sec:lockean-threshold-phase}

Gate~6 isolated an exact local condition for classical Lockean belief, but its
consistency component was still imposed semantically.  This section derives that
component from the stronger finite-probability layer and then separates the role of
the threshold itself from the legacy model contract.

The result is a sharper boundary.  In an integrity-certified model with a classical
accessible profile, positive and negative support are not merely two unrelated event
lists.  They form a complementary partition of the accessibility support.  The
legacy supermajority threshold \(c>1/2\) therefore excludes belief gluts
structurally.  The only remaining failure mode for classical belief is indecision.

\subsection{Classical profiles induce complementary support masses}

Fix an agent \(i\), a world \(w\), and an FDE-valued profile \(v\) on the worlds
accessible from \(w\).  Let
\[
  E^+ = \{u\in R_i(w) : v(u).\mathsf{pos}=\mathsf{true}\},
  \qquad
  E^- = \{u\in R_i(w) : v(u).\mathsf{neg}=\mathsf{true}\}.
\]
If every accessible value is classical, then each accessible world lies in exactly
one of these two events.  Hence they are disjoint and cover the accessibility
support.

\begin{theorem}[Complementary classical support masses]
\label{thm:classical-support-masses-sum-one}
Let \(m\) satisfy \texttt{ModelProbabilityIntegrity}.  If every value in the
accessible profile is in \(\{\Tval,\Fval\}\), then
\[
  \mu_{i,w}(E^+) + \mu_{i,w}(E^-) = 1.
\]
\statusverified
\end{theorem}

The proof uses the finite-probability integrity contract rather than the lightweight
legacy probability interface.  Duplicate-freedom of the accessibility support is
transported through the two filters, the events are shown disjoint, finite
additivity combines their masses, and extensionality identifies their union with the
whole support.

This theorem closes the precise gap left open at the end of Gate~6: once the input
profile is classical, the two threshold masses are genuine complements.

\subsection{Supermajority eliminates the glut branch}

Write
\[
  p=P^+(v),\qquad n=P^-(v).
\]
For a classical accessible profile under probability integrity,
Theorem~\ref{thm:classical-support-masses-sum-one} gives \(p+n=1\).  The existing
4-PEL model contract also requires
\[
  c>\frac12.
\]
These two facts are enough to rule out simultaneous threshold acceptance.

\begin{theorem}[Supermajority threshold consistency]
\label{thm:supermajority-threshold-consistency}
If \(p+n=1\) and \(c>1/2\), then
\[
  \neg(c\le p\land c\le n).
\]
Consequently, every integrity-certified classical accessible profile satisfies
\texttt{BeliefThresholdConsistent}.
\statusverified
\end{theorem}

Thus the exact Gate~6 recovery boundary simplifies in the stronger model class.

\begin{corollary}[Classical belief iff decisive belief]
\label{cor:belief-classical-iff-complete-integrity}
Under \texttt{ModelProbabilityIntegrity} and a classical accessible profile,
\[
  \Bel(v)\in\{\Tval,\Fval\}
  \quad\Longleftrightarrow\quad
  \mathsf{BelComplete}(v).
\]
Equivalently,
\[
  \Bel(v)\notin\{\Tval,\Fval\}
  \quad\Longleftrightarrow\quad
  \neg\mathsf{BelComplete}(v).
\]
\statusverified
\end{corollary}

The second form is conceptually useful.  Once the profile is classical and the
probability laws are certified, nonclassical Lockean belief is exactly a gap caused
by indecision.  A \(\Bval\)-valued belief is no longer a live possibility in this
sector.

\subsection{The model-independent threshold boundary}

The preceding theorem uses the existing model's built-in supermajority assumption.
To see what role \(1/2\) plays independently of that design choice, Gate~7 also
factors the threshold arithmetic into a model-free layer.  For rational masses
\(p,n\) with
\[
  p+n=1,
\]
define threshold completeness by \(c\le p\lor c\le n\), and threshold consistency
by \(\neg(c\le p\land c\le n)\).

Two complementary one-sided results follow.

\begin{theorem}[Threshold phase boundary]
\label{thm:threshold-phase-boundary}
For complementary masses \(p+n=1\):
\begin{enumerate}[label=(\roman*)]
  \item if \(c\le 1/2\), threshold completeness is guaranteed, so a gap is
        impossible;
  \item if \(c>1/2\), threshold consistency is guaranteed, so a glut is
        impossible.
\end{enumerate}
\statusverified
\end{theorem}

The critical value is therefore not merely a conventional midpoint.  It separates
two qualitatively different failure guarantees.  Lower thresholds protect against
indecision but permit overacceptance; strict supermajority thresholds protect
against overacceptance but permit indecision.

The inclusivity of the existing comparison matters at the boundary.  Because belief
uses \(\ge c\), an exact tie behaves as follows.

\begin{theorem}[Exact tie dichotomy]
\label{thm:half-half-phase-dichotomy}
For \(p=n=1/2\):
\[
  c\le\frac12
  \quad\Longrightarrow\quad
  c\le p\land c\le n,
\]
whereas
\[
  c>\frac12
  \quad\Longrightarrow\quad
  \neg(c\le p)\land\neg(c\le n).
\]
In particular, at \(c=1/2\) the tie is glut-like rather than gap-like.
\statusverified
\end{theorem}

In four-valued language, the threshold image of the tie jumps across the critical
value:
\[
  \boxed{
  c\le\frac12 : \Bval
  \qquad\vert\qquad
  c>\frac12 : \Nval
  }
\]
where the labels describe the two threshold bits, not the underlying world-level
truth value.

\subsection{No universal inclusive Lockean cutoff}

The tie immediately yields a sharp impossibility statement.

\begin{theorem}[No universal inclusive threshold]
\label{thm:no-universal-inclusive-threshold}
For every rational threshold \(c\), the complementary profile
\[
  p=n=\frac12
\]
is not threshold-regular.  Hence no single inclusive Lockean cutoff makes every
complementary probability profile classically two-valued.
\statusverified
\end{theorem}

This result should not be read as a defect of four-valued semantics.  It is a
property of inclusive thresholding on a complementary probability line.  A single
cutoff cannot simultaneously guarantee both completeness and consistency at the
exact midpoint.  The legacy 4-PEL choice \(c>1/2\) resolves the ambiguity in favor
of consistency, sending the balanced case to the gap side rather than the glut
side.

\subsection{Formula recovery with one fewer obligation}

Gate~6 defined \texttt{Formula.ClassicalRecoveryAdmissibleAt}; every belief node
required both threshold completeness and threshold consistency.  Gate~7 can now
weaken this recursive contract for integrity-certified models.

Define \texttt{Formula.ProbabilityRecoveryAdmissibleAt} exactly as before for atoms,
negation, and conjunction.  At a belief node, require only:
\begin{enumerate}[label=(\roman*)]
  \item recursive admissibility of the subformula throughout the accessible range;
  \item local threshold completeness.
\end{enumerate}
No independent threshold-consistency premise is included.

\begin{theorem}[Formula recovery under probability integrity]
\label{thm:probability-integrity-formula-recovery}
Let \(m\) satisfy \texttt{ModelProbabilityIntegrity}.  For every legacy formula
\(\varphi\),
\[
  \mathsf{ProbabilityRecoveryAdmissibleAt}(m,w,\varphi)
  \quad\Longrightarrow\quad
  \llbracket\varphi\rrbracket_{m,w}\in\{\Tval,\Fval\}.
\]
The corresponding global admissibility predicate yields classical evaluation at
every world.
\statusverified
\end{theorem}

At belief nodes the induction first obtains a classical accessible subformula
profile.  Probability integrity then supplies complementary masses, the built-in
supermajority threshold supplies consistency, and the sole explicit belief premise,
completeness, finishes the classical-recovery step.

This gives the Gate~7 hierarchy
\[
  \text{classical accessible profile}
  + \text{probability integrity}
  + c>\tfrac12
  \Longrightarrow
  \text{threshold consistency},
\]
followed by
\[
  \text{threshold completeness}
  \Longleftrightarrow
  \text{classical Lockean belief}.
\]

\paragraph{Gate boundary.}
Gate~7 derives the consistency half of Lockean classical recovery from the stronger
finite-probability semantics, formalizes the midpoint phase transition independently
of the model, proves the exact tie obstruction for every inclusive threshold, and
lifts the reduced recovery contract through the full legacy formula syntax.  It does
not prove that arbitrary formulas are threshold-complete, nor does it prove the
commented CPEL proof-system or translation completeness claim.  Decisiveness remains
a genuine semantic obligation.
